% MIT 18.03 Spring 2010 PS1--PS5；源和逐叶覆盖见 research/ps01-05-coverage.json。
\originalchapter{作业一：一阶方程}
\label{cw:ps1}
本章保留官方作业的讲次编号，因而从第0题开始。第一部分中的 EP 指 Edwards--Penney 教材；官方作业只给出节号与题号，没有给出书题题面。本书对这些条目只列定位，未翻译、未解答，也不把它们计入已收录题面。Notes 指课程公开的习题集，下面的链接回指本书相应中文题目与解答；字母表示只指定该题的相应小问。第二部分的题面来自官方作业，解答依据官方答案审译，并在必要处补证或修正。涉及 Mathlet 的图均按方程独立重绘，图示不是实际运行交互程序的记录。
\originalsection{第一部分：指定练习}
\begin{exercise}\label{cw:ps1-I-0}
原题 I.0（自然增长与可分离方程）：Notes 1A-5(c)，见\hyperref[cw:bank-1A-5]{题1A-5(c)}；EP 第1.1节第32、33、35题，第1.4节第39、66题，第1.5节第1、9、20题。
\end{exercise}
\begin{solution}
公开 Notes 题的中文题面与完整解答见上述链接。EP 八道题只有书题定位，本章没有其合法可用的完整题面，因此不代拟题目、不声称已翻译或解出这些书题。
\end{solution}
\begin{exercise}\label{cw:ps1-I-1}
原题 I.1（方向场与等倾线）：Notes 1C-1(a)、(b)、(e)，见\hyperref[cw:bank-1C-1]{题1C-1}。
\end{exercise}
\begin{solution}按所列三个小问阅读链接中的题面与解答；(c)、(d)并非本次作业指定。\end{solution}
\begin{exercise}\label{cw:ps1-I-2}
原题 I.2（欧拉法）：Notes 1C-4，见\hyperref[cw:bank-1C-4]{题1C-4}。
\end{exercise}
\begin{solution}题面、欧拉迭代表与计算见链接。\end{solution}
\begin{exercise}\label{cw:ps1-I-3}
原题 I.3（线性模型）：EP 第1.5节第33、45题。两题均须明确写出微分方程。
\end{exercise}
\begin{solution}仅列教材定位；书题题面未公开收入本次材料，未翻译、未解答。\end{solution}
\originalsection{第二部分：题面与解答}
\begin{exercise}\label{cw:ps1-II-0}
原题 II.0。羚羊种群在短时间 $\Delta t$ 内的变化近似为 $k(t)x(t)\Delta t$。不考虑捕猎，时间以年计，种群以千只羚羊（ko）计。肯尼亚塔纳河地区出现一种病毒，使自然增长率变为
\[
 k(t)=\frac{k_0}{(a+t)^2},\qquad t\geq0,\quad a,k_0>0.
\]
(a) 常数 $a,k_0$ 的单位是什么？(b) 写出模型方程。(c) 求通解，考虑数学上所有允许的 $t,x$，正确处理 $\ln|x|$ 和分离变量时丢失的解。(d) 若 $x(0)=x_0>0$，种群是否在大时间趋于稳定？若是，求极限。
\end{exercise}
\begin{solution}
解答依据（编者译审或补解，AI 辅助）：官方 PS1 解答第1页，补充定义区间。由 $x'=kx$ 可知 $k$ 的单位是年$^{-1}$，所以 $a$ 的单位是年，而 $k_0$ 的单位是年，保证指数无量纲。短时间增量关系除以 $\Delta t$ 后取极限得
\[
 x'=\frac{k_0x}{(a+t)^2}.
\]
在 $x\ne0$ 的区间分离变量，积分得 $\ln|x|=-k_0/(a+t)+c$，于是 $x=Ce^{-k_0/(a+t)}$，其中 $C$ 可正可负；把 $C=0$ 纳入便补回零解。方程在 $t=-a$ 无定义，通解应分别理解为 $(-\infty,-a)$ 或 $(-a,\infty)$ 内的解，不能跨越奇点。在种群模型的 $t\geq0$ 上，由初值得
\[
 x(t)=x_0\exp\left(\frac{k_0}{a}-\frac{k_0}{a+t}\right),
 \qquad\lim_{t\to\infty}x(t)=x_0e^{k_0/a}.
\]
增长率虽始终为正，其从0到无穷的积分却有限，所以种群增长到有限极限。
\end{solution}
\begin{exercise}\label{cw:ps1-II-1}
原题 II.1。研究 $y'=y^2-x$。
(a) 在横、纵轴均标出 $-4$ 至 $4$ 的大图上，画斜率 $m=-1,0,1,2$ 的等倾线，并在同一图上依据它们描出若干解。原作业随后要求用 Isoclines Mathlet 选此方程，再画 $m=-2,0,1,2$ 的等倾线，比较手绘图并拖动初值观察解。
(b) 描出分界解：其上方的解向右增长无界，其下方的解最终下降。
(c) 解曲线在 $(a,b)$ 与 $m=-1$ 的等倾线相切，求 $(a,b)$。
(d) 若 $y(a)<b$，给出解在 $x\to\infty$ 时所渐近的显式函数 $f(x)$，判断解最终位于其上还是其下。原题要求说明：(i) 解不能在 $x>a$ 穿越 $m=-1$ 的等倾线；(ii) 若在 $x=a$ 低于 $m=1$ 的等倾线，最终必穿越它；(iii) 解最终必穿越零倾线。
(e) 若 $(c,d)$ 是临界点，$c,d$ 满足什么关系？(f) 所有临界点都是局部极大值吗？须验证。
\end{exercise}
\begin{solution}
解答依据（编者译审或补解，AI 辅助）：官方 PS1 解答第1--2页；(d)补充渐近收敛证明，并说明原题(iii)的条件缺漏。
等倾线为 $x=y^2-m$。沿线画斜率为 $m$ 的短线段便得到方向场；零倾线 $x=y^2$ 外侧 $y^2>x$ 的区域斜率为正，内侧斜率为负。下图直接由这些等式重绘，各条线段只表示局部斜率；灰色解曲线由 $y(0)=0.5,1.4$ 的数值积分叠加，供方向场描线的比较。
\begin{center}
\begin{tikzpicture}
\begin{axis}[width=.78\linewidth,height=6cm,xmin=-4,xmax=4,ymin=-4,ymax=4,domain=-4:4,axis lines=middle,xlabel={$x$},ylabel={$y$},samples=100,legend style={at={(1.02,1)},anchor=north west,font=\small}]
\addplot[domain=-2:2,blue] ({x*x+1},x);\addlegendentry{$m=-1$}
\addplot[domain=-2:2,black] ({x*x},x);\addlegendentry{$m=0$}
\addplot[domain=-2.2:2.2,red] ({x*x-1},x);\addlegendentry{$m=1$}
\addplot[domain=-2.4:2.4,orange] ({x*x-2},x);\addlegendentry{$m=2$}
\addplot[gray,thick,forget plot]coordinates{(0.000000,0.500000)(0.054054,0.512401)(0.108108,0.522503)(0.162162,0.530185)(0.216216,0.535313)(0.270270,0.537738)(0.324324,0.537300)(0.378378,0.533828)(0.432432,0.527141)(0.486486,0.517057)(0.540541,0.503389)(0.594595,0.485955)(0.648649,0.464583)(0.702703,0.439116)(0.756757,0.409422)(0.810811,0.375399)(0.864865,0.336990)(0.918919,0.294184)(0.972973,0.247030)(1.027027,0.195644)(1.081081,0.140213)(1.135135,0.080995)(1.189189,0.018328)(1.243243,-0.047383)(1.297297,-0.115667)(1.351351,-0.186003)(1.405405,-0.257827)(1.459459,-0.330555)(1.513514,-0.403598)(1.567568,-0.476380)(1.621622,-0.548356)(1.675676,-0.619026)(1.729730,-0.687947)(1.783784,-0.754746)(1.837838,-0.819117)(1.891892,-0.880831)(1.945946,-0.939727)(2.000000,-0.995714)(2.054054,-1.048758)(2.108108,-1.098881)(2.162162,-1.146145)(2.216216,-1.190650)(2.270270,-1.232522)(2.324324,-1.271907)(2.378378,-1.308963)(2.432432,-1.343854)(2.486486,-1.376747)(2.540541,-1.407806)(2.594595,-1.437191)(2.648649,-1.465053)(2.702703,-1.491536)(2.756757,-1.516774)(2.810811,-1.540889)(2.864865,-1.563996)(2.918919,-1.586196)(2.972973,-1.607582)(3.027027,-1.628239)(3.081081,-1.648240)(3.135135,-1.667653)(3.189189,-1.686536)(3.243243,-1.704943)(3.297297,-1.722920)(3.351351,-1.740507)(3.405405,-1.757740)(3.459459,-1.774652)(3.513514,-1.791270)(3.567568,-1.807618)(3.621622,-1.823718)(3.675676,-1.839589)(3.729730,-1.855246)(3.783784,-1.870706)(3.837838,-1.885980)(3.891892,-1.901079)(3.945946,-1.916014)(4.000000,-1.930794)};
\addplot[gray,thick,forget plot]coordinates{(0.000000,1.400000)(0.009970,1.419769)(0.019941,1.440004)(0.029911,1.460728)(0.039882,1.481964)(0.049852,1.503735)(0.059823,1.526069)(0.069793,1.548992)(0.079763,1.572533)(0.089734,1.596723)(0.099704,1.621595)(0.109675,1.647183)(0.119645,1.673524)(0.129616,1.700658)(0.139586,1.728628)(0.149556,1.757477)(0.159527,1.787254)(0.169497,1.818010)(0.179468,1.849801)(0.189438,1.882684)(0.199409,1.916725)(0.209379,1.951991)(0.219349,1.988556)(0.229320,2.026498)(0.239290,2.065903)(0.249261,2.106865)(0.259231,2.149483)(0.269202,2.193866)(0.279172,2.240132)(0.289143,2.288411)(0.299113,2.338842)(0.309083,2.391580)(0.319054,2.446792)(0.329024,2.504664)(0.338995,2.565398)(0.348965,2.629219)(0.358936,2.696373)(0.368906,2.767136)(0.378876,2.841812)(0.388847,2.920741)(0.398817,3.004302)(0.408788,3.092921)(0.418758,3.187077)(0.428729,3.287311)(0.438699,3.394234)(0.448669,3.508546)(0.458640,3.631041)(0.468610,3.762636)(0.478581,3.904386)};

\foreach \xx in {-3,-2,-1,0,1,2,3}{\foreach \yy in {-3,-2,-1,0,1,2,3}{\pgfmathsetmacro{\mm}{\yy*\yy-\xx}\pgfmathsetmacro{\dx}{.18/sqrt(1+\mm*\mm)}\pgfmathsetmacro{\dy}{\mm*\dx}\edef\fieldsegment{\noexpand\draw[gray] (axis cs:{\xx-\dx},{\yy-\dy})--(axis cs:{\xx+\dx},{\yy+\dy});}\fieldsegment}}
\end{axis}
\end{tikzpicture}
\end{center}
(b) 分界解的数值示意另见下图；它向右沿正支 $\sqrt x$ 上升，而较低的解转向负支。数值图用于定位分界，不作为所有解分类的证明。可用代换 $y=-u'/u$ 得到线性方程 $u''=xu$；取向右衰减的非零解 $u$ 所得到的曲线正是此分界解。因题目要求描图，下面采用从 $x=16$ 的近似渐近初值向左积分的曲线，并明确保留其数值性质。
% Numerical coordinates are produced from the equation, not from the Mathlet.
\begin{center}\begin{tikzpicture}\begin{axis}[width=.75\linewidth,height=5cm,axis lines=middle,xlabel={$x$},ylabel={$y$},xmin=-3,xmax=4,ymin=-3,ymax=4,domain=-3:4,samples=100]
\addplot[blue,thick]coordinates{(-2.026816,-2.990891)(-1.992179,-2.646832)(-1.957542,-2.361705)(-1.922905,-2.120841)(-1.888268,-1.914086)(-1.853631,-1.734187)(-1.818994,-1.575824)(-1.784358,-1.435004)(-1.749721,-1.308673)(-1.715084,-1.194455)(-1.680447,-1.090471)(-1.645810,-0.995220)(-1.611173,-0.907481)(-1.576536,-0.826258)(-1.541899,-0.750725)(-1.507263,-0.680193)(-1.472626,-0.614084)(-1.437989,-0.551905)(-1.403352,-0.493239)(-1.368715,-0.437726)(-1.334078,-0.385054)(-1.299441,-0.334954)(-1.264804,-0.287191)(-1.230168,-0.241558)(-1.195531,-0.197874)(-1.160894,-0.155977)(-1.126257,-0.115724)(-1.091620,-0.076989)(-1.056983,-0.039657)(-1.022346,-0.003627)(-0.987709,0.031195)(-0.953073,0.064890)(-0.918436,0.097533)(-0.883799,0.129194)(-0.849162,0.159933)(-0.814525,0.189808)(-0.779888,0.218871)(-0.745251,0.247168)(-0.710615,0.274743)(-0.675978,0.301636)(-0.641341,0.327885)(-0.606704,0.353523)(-0.572067,0.378582)(-0.537430,0.403090)(-0.502793,0.427076)(-0.468156,0.450564)(-0.433520,0.473577)(-0.398883,0.496139)(-0.364246,0.518268)(-0.329609,0.539985)(-0.294972,0.561306)(-0.260335,0.582249)(-0.225698,0.602830)(-0.191061,0.623064)(-0.156425,0.642963)(-0.121788,0.662542)(-0.087151,0.681812)(-0.052514,0.700786)(-0.017877,0.719474)(0.016760,0.737887)(0.051397,0.756034)(0.086034,0.773925)(0.120670,0.791569)(0.155307,0.808975)(0.189944,0.826149)(0.224581,0.843100)(0.259218,0.859835)(0.293855,0.876360)(0.328492,0.892683)(0.363128,0.908809)(0.397765,0.924745)(0.432402,0.940496)(0.467039,0.956068)(0.501676,0.971465)(0.536313,0.986693)(0.570950,1.001757)(0.605587,1.016660)(0.640223,1.031408)(0.674860,1.046004)(0.709497,1.060453)(0.744134,1.074759)(0.778771,1.088924)(0.813408,1.102953)(0.848045,1.116849)(0.882682,1.130616)(0.917318,1.144256)(0.951955,1.157772)(0.986592,1.171168)(1.021229,1.184446)(1.055866,1.197610)(1.090503,1.210661)(1.125140,1.223602)(1.159777,1.236436)(1.194413,1.249164)(1.229050,1.261790)(1.263687,1.274316)(1.298324,1.286743)(1.332961,1.299074)(1.367598,1.311311)(1.402235,1.323455)(1.436872,1.335508)(1.471508,1.347473)(1.506145,1.359351)(1.540782,1.371144)(1.575419,1.382853)(1.610056,1.394480)(1.644693,1.406027)(1.679330,1.417495)(1.713966,1.428885)(1.748603,1.440200)(1.783240,1.451439)(1.817877,1.462606)(1.852514,1.473700)(1.887151,1.484724)(1.921788,1.495678)(1.956425,1.506564)(1.991061,1.517382)(2.025698,1.528135)(2.060335,1.538823)(2.094972,1.549447)(2.129609,1.560008)(2.164246,1.570507)(2.198883,1.580946)(2.233520,1.591325)(2.268156,1.601645)(2.302793,1.611907)(2.337430,1.622112)(2.372067,1.632261)(2.406704,1.642355)(2.441341,1.652394)(2.475978,1.662380)(2.510615,1.672313)(2.545251,1.682194)(2.579888,1.692023)(2.614525,1.701802)(2.649162,1.711531)(2.683799,1.721211)(2.718436,1.730843)(2.753073,1.740426)(2.787709,1.749963)(2.822346,1.759454)(2.856983,1.768898)(2.891620,1.778297)(2.926257,1.787652)(2.960894,1.796963)(2.995531,1.806230)(3.030168,1.815455)(3.064804,1.824637)(3.099441,1.833778)(3.134078,1.842877)(3.168715,1.851936)(3.203352,1.860954)(3.237989,1.869933)(3.272626,1.878873)(3.307263,1.887774)(3.341899,1.896637)(3.376536,1.905462)(3.411173,1.914250)(3.445810,1.923001)(3.480447,1.931715)(3.515084,1.940394)(3.549721,1.949037)(3.584358,1.957645)(3.618994,1.966218)(3.653631,1.974757)(3.688268,1.983261)(3.722905,1.991733)(3.757542,2.000171)(3.792179,2.008576)(3.826816,2.016949)(3.861453,2.025290)(3.896089,2.033599)(3.930726,2.041877)(3.965363,2.050123)(4.000000,2.058339)};
\addplot[gray]coordinates{(0.000000,-1.000000)(0.036697,-0.965259)(0.073394,-0.934197)(0.110092,-0.906495)(0.146789,-0.881881)(0.183486,-0.860114)(0.220183,-0.840985)(0.256881,-0.824307)(0.293578,-0.809916)(0.330275,-0.797664)(0.366972,-0.787416)(0.403670,-0.779053)(0.440367,-0.772463)(0.477064,-0.767546)(0.513761,-0.764208)(0.550459,-0.762362)(0.587156,-0.761926)(0.623853,-0.762823)(0.660550,-0.764982)(0.697248,-0.768331)(0.733945,-0.772807)(0.770642,-0.778345)(0.807339,-0.784884)(0.844037,-0.792366)(0.880734,-0.800733)(0.917431,-0.809930)(0.954128,-0.819904)(0.990826,-0.830602)(1.027523,-0.841973)(1.064220,-0.853969)(1.100917,-0.866542)(1.137615,-0.879644)(1.174312,-0.893231)(1.211009,-0.907259)(1.247706,-0.921686)(1.284404,-0.936471)(1.321101,-0.951576)(1.357798,-0.966962)(1.394495,-0.982594)(1.431193,-0.998438)(1.467890,-1.014460)(1.504587,-1.030630)(1.541284,-1.046919)(1.577982,-1.063300)(1.614679,-1.079746)(1.651376,-1.096234)(1.688073,-1.112740)(1.724771,-1.129246)(1.761468,-1.145731)(1.798165,-1.162177)(1.834862,-1.178570)(1.871560,-1.194895)(1.908257,-1.211138)(1.944954,-1.227288)(1.981651,-1.243334)(2.018349,-1.259268)(2.055046,-1.275081)(2.091743,-1.290768)(2.128440,-1.306321)(2.165138,-1.321736)(2.201835,-1.337010)(2.238532,-1.352138)(2.275229,-1.367119)(2.311927,-1.381952)(2.348624,-1.396634)(2.385321,-1.411165)(2.422018,-1.425546)(2.458716,-1.439777)(2.495413,-1.453859)(2.532110,-1.467793)(2.568807,-1.481580)(2.605505,-1.495223)(2.642202,-1.508723)(2.678899,-1.522083)(2.715596,-1.535305)(2.752294,-1.548391)(2.788991,-1.561345)(2.825688,-1.574168)(2.862385,-1.586864)(2.899083,-1.599435)(2.935780,-1.611884)(2.972477,-1.624214)(3.009174,-1.636428)(3.045872,-1.648529)(3.082569,-1.660519)(3.119266,-1.672401)(3.155963,-1.684178)(3.192661,-1.695852)(3.229358,-1.707426)(3.266055,-1.718903)(3.302752,-1.730284)(3.339450,-1.741573)(3.376147,-1.752771)(3.412844,-1.763881)(3.449541,-1.774905)(3.486239,-1.785845)(3.522936,-1.796704)(3.559633,-1.807482)(3.596330,-1.818183)(3.633028,-1.828807)(3.669725,-1.839357)(3.706422,-1.849834)(3.743119,-1.860241)(3.779817,-1.870578)(3.816514,-1.880847)(3.853211,-1.891050)(3.889908,-1.901188)(3.926606,-1.911263)(3.963303,-1.921275)(4.000000,-1.931227)};
\addplot[gray]coordinates{(0.000000,0.000000)(0.036697,-0.000673)(0.073394,-0.002693)(0.110092,-0.006059)(0.146789,-0.010770)(0.183486,-0.016823)(0.220183,-0.024215)(0.256881,-0.032938)(0.293578,-0.042985)(0.330275,-0.054345)(0.366972,-0.067004)(0.403670,-0.080943)(0.440367,-0.096142)(0.477064,-0.112576)(0.513761,-0.130216)(0.550459,-0.149027)(0.587156,-0.168973)(0.623853,-0.190011)(0.660550,-0.212094)(0.697248,-0.235172)(0.733945,-0.259190)(0.770642,-0.284089)(0.807339,-0.309806)(0.844037,-0.336276)(0.880734,-0.363431)(0.917431,-0.391199)(0.954128,-0.419509)(0.990826,-0.448286)(1.027523,-0.477456)(1.064220,-0.506944)(1.100917,-0.536677)(1.137615,-0.566582)(1.174312,-0.596588)(1.211009,-0.626626)(1.247706,-0.656629)(1.284404,-0.686535)(1.321101,-0.716285)(1.357798,-0.745823)(1.394495,-0.775098)(1.431193,-0.804062)(1.467890,-0.832675)(1.504587,-0.860898)(1.541284,-0.888697)(1.577982,-0.916045)(1.614679,-0.942917)(1.651376,-0.969293)(1.688073,-0.995158)(1.724771,-1.020500)(1.761468,-1.045310)(1.798165,-1.069585)(1.834862,-1.093321)(1.871560,-1.116522)(1.908257,-1.139190)(1.944954,-1.161332)(1.981651,-1.182956)(2.018349,-1.204071)(2.055046,-1.224689)(2.091743,-1.244822)(2.128440,-1.264485)(2.165138,-1.283691)(2.201835,-1.302455)(2.238532,-1.320793)(2.275229,-1.338720)(2.311927,-1.356252)(2.348624,-1.373405)(2.385321,-1.390194)(2.422018,-1.406634)(2.458716,-1.422742)(2.495413,-1.438531)(2.532110,-1.454016)(2.568807,-1.469210)(2.605505,-1.484128)(2.642202,-1.498783)(2.678899,-1.513186)(2.715596,-1.527349)(2.752294,-1.541285)(2.788991,-1.555003)(2.825688,-1.568514)(2.862385,-1.581828)(2.899083,-1.594954)(2.935780,-1.607901)(2.972477,-1.620677)(3.009174,-1.633290)(3.045872,-1.645746)(3.082569,-1.658054)(3.119266,-1.670220)(3.155963,-1.682250)(3.192661,-1.694149)(3.229358,-1.705923)(3.266055,-1.717577)(3.302752,-1.729116)(3.339450,-1.740544)(3.376147,-1.751866)(3.412844,-1.763086)(3.449541,-1.774207)(3.486239,-1.785232)(3.522936,-1.796166)(3.559633,-1.807011)(3.596330,-1.817770)(3.633028,-1.828446)(3.669725,-1.839042)(3.706422,-1.849559)(3.743119,-1.860001)(3.779817,-1.870368)(3.816514,-1.880665)(3.853211,-1.890891)(3.889908,-1.901050)(3.926606,-1.911143)(3.963303,-1.921171)(4.000000,-1.931136)};
\addplot[gray]coordinates{(0.000000,0.500000)(0.036697,0.508664)(0.073394,0.516285)(0.110092,0.522828)(0.146789,0.528254)(0.183486,0.532521)(0.220183,0.535585)(0.256881,0.537398)(0.293578,0.537910)(0.330275,0.537070)(0.366972,0.534822)(0.403670,0.531110)(0.440367,0.525878)(0.477064,0.519068)(0.513761,0.510620)(0.550459,0.500478)(0.587156,0.488585)(0.623853,0.474886)(0.660550,0.459333)(0.697248,0.441876)(0.733945,0.422476)(0.770642,0.401098)(0.807339,0.377716)(0.844037,0.352311)(0.880734,0.324878)(0.917431,0.295420)(0.954128,0.263957)(0.990826,0.230519)(1.027523,0.195154)(1.064220,0.157923)(1.100917,0.118905)(1.137615,0.078194)(1.174312,0.035898)(1.211009,-0.007856)(1.247706,-0.052930)(1.284404,-0.099173)(1.321101,-0.146420)(1.357798,-0.194502)(1.394495,-0.243238)(1.431193,-0.292446)(1.467890,-0.341941)(1.504587,-0.391539)(1.541284,-0.441059)(1.577982,-0.490326)(1.614679,-0.539175)(1.651376,-0.587448)(1.688073,-0.635003)(1.724771,-0.681707)(1.761468,-0.727447)(1.798165,-0.772120)(1.834862,-0.815641)(1.871560,-0.857941)(1.908257,-0.898964)(1.944954,-0.938672)(1.981651,-0.977039)(2.018349,-1.014050)(2.055046,-1.049704)(2.091743,-1.084012)(2.128440,-1.116990)(2.165138,-1.148666)(2.201835,-1.179072)(2.238532,-1.208248)(2.275229,-1.236237)(2.311927,-1.263086)(2.348624,-1.288844)(2.385321,-1.313562)(2.422018,-1.337293)(2.458716,-1.360087)(2.495413,-1.381998)(2.532110,-1.403076)(2.568807,-1.423371)(2.605505,-1.442932)(2.642202,-1.461805)(2.678899,-1.480035)(2.715596,-1.497664)(2.752294,-1.514734)(2.788991,-1.531282)(2.825688,-1.547345)(2.862385,-1.562956)(2.899083,-1.578148)(2.935780,-1.592950)(2.972477,-1.607389)(3.009174,-1.621492)(3.045872,-1.635282)(3.082569,-1.648782)(3.119266,-1.662011)(3.155963,-1.674989)(3.192661,-1.687732)(3.229358,-1.700258)(3.266055,-1.712580)(3.302752,-1.724712)(3.339450,-1.736666)(3.376147,-1.748454)(3.412844,-1.760086)(3.449541,-1.771572)(3.486239,-1.782920)(3.522936,-1.794138)(3.559633,-1.805235)(3.596330,-1.816215)(3.633028,-1.827086)(3.669725,-1.837852)(3.706422,-1.848520)(3.743119,-1.859094)(3.779817,-1.869578)(3.816514,-1.879976)(3.853211,-1.890291)(3.889908,-1.900528)(3.926606,-1.910689)(3.963303,-1.920776)(4.000000,-1.930794)};
\addplot[gray]coordinates{(0.000000,1.000000)(0.010159,1.010211)(0.020318,1.020530)(0.030477,1.030961)(0.040636,1.041508)(0.050795,1.052176)(0.060954,1.062971)(0.071113,1.073897)(0.081272,1.084960)(0.091431,1.096165)(0.101590,1.107518)(0.111749,1.119025)(0.121908,1.130692)(0.132067,1.142526)(0.142226,1.154533)(0.152385,1.166721)(0.162544,1.179097)(0.172703,1.191669)(0.182862,1.204444)(0.193021,1.217431)(0.203180,1.230639)(0.213339,1.244077)(0.223498,1.257755)(0.233657,1.271682)(0.243816,1.285868)(0.253975,1.300326)(0.264134,1.315067)(0.274293,1.330102)(0.284452,1.345444)(0.294611,1.361107)(0.304770,1.377104)(0.314929,1.393451)(0.325088,1.410162)(0.335247,1.427255)(0.345406,1.444745)(0.355565,1.462652)(0.365724,1.480995)(0.375883,1.499793)(0.386042,1.519068)(0.396200,1.538842)(0.406359,1.559140)(0.416518,1.579986)(0.426677,1.601407)(0.436836,1.623432)(0.446995,1.646090)(0.457154,1.669415)(0.467313,1.693439)(0.477472,1.718199)(0.487631,1.743734)(0.497790,1.770084)(0.507949,1.797295)(0.518108,1.825413)(0.528267,1.854488)(0.538426,1.884574)(0.548585,1.915729)(0.558744,1.948017)(0.568903,1.981502)(0.579062,2.016258)(0.589221,2.052362)(0.599380,2.089898)(0.609539,2.128958)(0.619698,2.169638)(0.629857,2.212047)(0.640016,2.256300)(0.650175,2.302524)(0.660334,2.350856)(0.670493,2.401448)(0.680652,2.454464)(0.690811,2.510085)(0.700970,2.568511)(0.711129,2.629962)(0.721288,2.694681)(0.731447,2.762937)(0.741606,2.835028)(0.751765,2.911289)(0.761924,2.992092)(0.772083,3.077855)(0.782242,3.169048)(0.792401,3.266200)(0.802560,3.369914)(0.812719,3.480875)(0.822878,3.599863)(0.833037,3.727777)(0.843196,3.865653)};
\addplot[gray]coordinates{(0.000000,1.400000)(0.006769,1.413371)(0.013538,1.426954)(0.020307,1.440756)(0.027076,1.454783)(0.033845,1.469043)(0.040614,1.483543)(0.047382,1.498291)(0.054151,1.513295)(0.060920,1.528563)(0.067689,1.544105)(0.074458,1.559928)(0.081227,1.576042)(0.087996,1.592458)(0.094765,1.609185)(0.101534,1.626235)(0.108303,1.643617)(0.115072,1.661345)(0.121841,1.679429)(0.128609,1.697883)(0.135378,1.716719)(0.142147,1.735953)(0.148916,1.755597)(0.155685,1.775667)(0.162454,1.796179)(0.169223,1.817150)(0.175992,1.838597)(0.182761,1.860537)(0.189530,1.882991)(0.196299,1.905979)(0.203068,1.929521)(0.209836,1.953639)(0.216605,1.978358)(0.223374,2.003701)(0.230143,2.029695)(0.236912,2.056366)(0.243681,2.083744)(0.250450,2.111859)(0.257219,2.140742)(0.263988,2.170429)(0.270757,2.200954)(0.277526,2.232357)(0.284295,2.264676)(0.291063,2.297955)(0.297832,2.332239)(0.304601,2.367576)(0.311370,2.404018)(0.318139,2.441619)(0.324908,2.480437)(0.331677,2.520534)(0.338446,2.561977)(0.345215,2.604835)(0.351984,2.649186)(0.358753,2.695109)(0.365522,2.742693)(0.372290,2.792030)(0.379059,2.843221)(0.385828,2.896374)(0.392597,2.951607)(0.399366,3.009044)(0.406135,3.068824)(0.412904,3.131092)(0.419673,3.196011)(0.426442,3.263753)(0.433211,3.334510)(0.439980,3.408487)(0.446749,3.485912)(0.453517,3.567032)(0.460286,3.652119)(0.467055,3.741473)(0.473824,3.835423)(0.480593,3.934334)};
\addplot[black,dashed,domain=0:4]{-sqrt(x)};\end{axis}\end{tikzpicture}\end{center}

(c) 在等倾线 $x=y^2+1$ 上，其切线满足 $1=2yy'$。相切时两条曲线斜率都为 $-1$，故 $b=-1/2$，$a=b^2+1=5/4$。

(d) 这里 $m=-1$ 线须取负支 $U(x)=-\sqrt{x-1}$。在 $x>a$，$U'(x)=-1/(2\sqrt{x-1})>-1$，而解在接触这条线时斜率为 $-1$。从下方第一次穿越必须使 $y'-U'\geq0$，与此矛盾，故 $y<U$，这证明(i)。解在任意有限的 $x$ 区间也不会逃向负无穷：若 $y$ 足够负，则 $y^2-x>0$；配合上界 $U$，得到可延拓到所有 $x\geq a$ 的解。

令 $L(x)=-\sqrt{x+1}$，即 $m=1$ 线的负支。若始终 $y<L$，则 $y'>1$，于是 $y(x)>y(a)+(x-a)$，最终与 $y<L<0$ 矛盾。因此必穿越 $L$，这证明(ii)。同理，若解还低于零倾线负支 $N(x)=-\sqrt x$，则 $y'>0$，而 $N\to-\infty$，故必穿越 $N$；在接触 $N$ 时 $y'-N'=1/(2\sqrt x)>0$，所以穿越后永远在其上。

原题(iii)未写出“起初低于零倾线”的条件，作为无条件陈述并不正确。例如取 $y(a)=-1$，此时 $-\sqrt{5/4}<-1<-1/2$，解从一开始就在 $N$ 上方，并永远不穿越它。正确结论是：每个满足本问初值的解最终都处于 $N$ 上方；初始位于其下方的解必穿越一次。最终
\[
 -\sqrt x<y(x)<-\sqrt{x-1},\qquad
 0<y(x)+\sqrt x<\frac1{\sqrt x+\sqrt{x-1}}\longrightarrow0.
\]
所以 $f(x)=-\sqrt x$，且最终从上方渐近。这一夹逼才证明曲线间的纵向距离趋于0。

(e) 在临界点 $0=y'(c)=d^2-c$，故 $c=d^2$。(f) 对方程求导得 $y''=2yy'-1$。临界点处 $y''=-1<0$，故每个临界点都是严格局部极大值；此结论不依赖图像猜测。
\end{solution}
\begin{exercise}\label{cw:ps1-II-2}
原题 II.2。(a) 设 $y'=2x$、$y(0)=0$。求 $y(1)$，并分别用2、3、$n$个等长步的欧拉法估计它。可用 $1+\cdots+(n-1)=n(n-1)/2$。估计是否收敛？(b) 步长 $h=1/n$ 时的精确误差是多少？是否符合误差至多正比于 $h$ 的预测？
\end{exercise}
\begin{solution}
解答依据（编者译审或补解，AI 辅助）：官方 PS1 解答第2--3页，修正其表头把斜率误写为 $-y_k$ 的笔误。积分给 $y=x^2$，所以 $y(1)=1$。令 $x_k=kh$，欧拉公式为 $y_{k+1}=y_k+2kh^2$，从 $y_0=0$ 得
\[
 y_n=2h^2\sum_{k=0}^{n-1}k=n(n-1)h^2=1-\frac1n.
\]
2步得到 $1/2$，3步得到 $2/3$，一般估计趋于1。真值减去估计的误差恰为 $1/n=h$，因此不仅符合一阶误差预测，还给出本题的确切误差常数1。
\end{solution}
\begin{exercise}\label{cw:ps1-II-3}
原题 II.3。史高治为侄子设立信托，投资按恒定年利率 $I$ 连续增长，每年恒定支付 $q$ 美元；$I$ 的单位为年$^{-1}$，例如 $I=0.05$ 表示每年5\%。(a) 建模并解释。(b) 求通解。(c) 取 $I=0.05$，若每月支付1000美元并永远保持本金不变，应投入多少？(d) 若仍每月支付1000美元，但要求恰好20年后用完，应投入多少？精确到美分。
\end{exercise}
\begin{solution}
解答依据（编者译审或补解，AI 辅助）：官方 PS1 解答第3页。设 $x(t)$ 为账户美元余额、$t$ 为年。在 $\Delta t$ 内利息约为 $Ix(t)\Delta t$、支出为 $q\Delta t$，故 $x'=Ix-q$。若 $I\ne0$，令 $u=x-q/I$，则 $u'=Iu$，所以
\[
 x(t)=\frac qI+Ce^{It};
\]
其中 $C=0$ 包含恒定余额解。若 $I=0$，则 $x=C-qt$。
(c) $q=12000$ 美元/年，恒定余额为 $q/I=240000$ 美元。
(d) 令 $T=20$，终值条件 $0=q/I+Ce^{IT}$ 给 $C=-(q/I)e^{-IT}$，所以
\[
 x(0)=\frac qI(1-e^{-IT})=240000(1-e^{-1})
 \simeq151708.93\ \text{美元}.
\]
按同一公式，$0\leq t<20$ 的余额为正，且在 $t=20$ 为零，符合模型要求。
\end{solution}

\originalchapter{作业二：复数与线性响应}\label{cw:ps2}
\originalsection{第一部分：指定练习}
\begin{exercise}\label{cw:ps2-I-4}原题 I.4：EP 第1.5节第1、2、5题；应先检查是否可分离变量。\end{exercise}
\begin{solution}仅列教材定位，未取得书题题面，未翻译、未解答；第1题也在 PS1 I.0 指定过。\end{solution}
\begin{exercise}\label{cw:ps2-I-5}原题 I.5：Notes 2E-1、2、7，分别见\hyperref[cw:bank-2E-1]{题2E-1}、\hyperref[cw:bank-2E-2]{题2E-2}、\hyperref[cw:bank-2E-7]{题2E-7}。\end{exercise}
\begin{solution}题面与解答复用所链公开题。\end{solution}
\begin{exercise}\label{cw:ps2-I-6}原题 I.6：(a) Notes 2E-9、10，见\hyperref[cw:bank-2E-9]{题2E-9}、\hyperref[cw:bank-2E-10]{题2E-10}。(b) 将下列函数写为 $A\cos(\omega t-\phi)$，每题先画对应直角三角形：(i) $\cos2t+\sin2t$；(ii) $\cos\pi t-\sqrt3\sin\pi t$；(iii) $\cos(t-\pi/8)+\sin(t-\pi/8)$。\end{exercise}
\begin{solution}
解答依据（编者译审或补解，AI 辅助）：PS2题面第3页答案；(iii)修正其把 $3\pi/8$ 排成 $3t/8$ 的笔误。(a)见链接。(b) 若余弦、正弦系数为 $(a,b)$，取斜边 $A=\sqrt{a^2+b^2}$，角 $\phi$ 满足 $\cos\phi=a/A$、$\sin\phi=b/A$。因此三题分别为
\[
 \sqrt2\cos(2t-\pi/4),\quad 2\cos(\pi t+\pi/3),\quad
 \sqrt2\cos(t-3\pi/8).
\]
对应三角形的有向直角边为 $(1,1)$、$(1,-\sqrt3)$、$(1,1)$；第三题先对自变量 $s=t-\pi/8$ 作同一个三角形，再将相位加回。
\begin{center}\begin{tikzpicture}[scale=.7]
\foreach \offset/\aa/\bb/\texta/\textb in {0/1.7/1.7/1/1,4/1.4/-2.42/1/{-\sqrt3},8/1.7/1.7/1/1}{\begin{scope}[xshift=\offset cm]\draw[->](-.2,0)--(2.4,0);\draw[->](0,-2.6)--(0,2.3);\draw[thick](0,0)--(\aa,0)--(\aa,\bb)--cycle;\node[below]at(\aa/2,0){$\texta$};\node[right]at(\aa,\bb/2){$\textb$};\end{scope}}
\end{tikzpicture}\end{center}
\end{solution}
\begin{exercise}\label{cw:ps2-I-7}原题 I.7：(a) Notes 2E-15，见\hyperref[cw:bank-2E-15]{题2E-15}。(b) 分别对 (i) $x'+2x=e^{3t}$，(ii) 复值方程 $x'+2x=e^{3it}$，先找与右端同指数形式的特解，再写通解。\end{exercise}
\begin{solution}解答依据（编者译审或补解，AI 辅助）：PS2题面第3页。(a)见链接。(b)代入 $Be^{\lambda t}$ 得 $(\lambda+2)B=1$，齐次解为 $Ce^{-2t}$，故分别为 $e^{3t}/5+Ce^{-2t}$（$C\in\R$）与 $e^{3it}/(2+3i)+Ce^{-2t}$（$C\in\C$）。\end{solution}
\originalsection{第二部分：题面与解答}
\begin{exercise}\label{cw:ps2-II-4}
原题 II.4。虚构元素 St 的半衰期为 $t_S$，每次衰变以相同概率产生 Mi 或稳定元素 En；Mi 的半衰期为 $t_M$，衰变为 En。初始只有 St，选单位使 $x(0)=1$；$x,y,z$ 分别表示 St、Mi、En 的量，$y(0)=z(0)=0$。题面明确 $t_M\ne t_S$（原 PDF 的不等号在纯文本提取中可能丢失斜线，须以原页为准）。
(a) 描图并求三者的长期极限。(b) 用衰变常数 $\sigma,\mu$ 建立方程，先联系半衰期，再检验总量守恒。(c) 依次求解。(d) Mi 何时最多？(e) 若 $x(0)=2$ 有何变化？(f) 无关的问题：若 $x=e^t$ 是 $tx'+2x=q(t)$ 的解，求 $q$ 及通解。
\end{exercise}
\begin{solution}
解答依据（编者译审，AI 辅助）：官方 PS2 解答第1页，补充定义区间与相等衰变率的极限对照；实际题目要求两半衰期不同。
由 $e^{-\sigma t_S}=1/2$ 得 $\sigma=(\ln2)/t_S$，同理 $\mu=(\ln2)/t_M$；两者都为正且 $\mu\ne\sigma$。每单位时间 St 减少 $\sigma x$，其中一半进入 Mi，因此
\[
 x'=-\sigma x,\quad y'=\tfrac12\sigma x-\mu y,
 \quad z'=\tfrac12\sigma x+\mu y.
\]
相加恰为零，初始总量1，故 $x+y+z=1$。先得 $x=e^{-\sigma t}$；乘 $y$ 方程以 $e^{\mu t}$ 得
\[
 (e^{\mu t}y)'=\tfrac12\sigma e^{(\mu-\sigma)t}.
\]
从0积分，利用 $y(0)=0$，可得
\[
 y=\frac\sigma{2(\mu-\sigma)}(e^{-\sigma t}-e^{-\mu t}),\qquad
 z=1-e^{-\sigma t}-\frac\sigma{2(\mu-\sigma)}(e^{-\sigma t}-e^{-\mu t}).
\]
$x$ 单调下降到0；$z'=\sigma x/2+\mu y\geq0$，从0上升到1。对 $y$ 求导得
\[
 y'=\frac\sigma{2(\mu-\sigma)}(-\sigma e^{-\sigma t}+\mu e^{-\mu t}).
\]
它在0处为 $\sigma/2>0$，只有一个零点：由 $\sigma e^{-\sigma t}=\mu e^{-\mu t}$ 得
\[
 t_{\max}=\frac{\ln(\mu/\sigma)}{\mu-\sigma}>0.
\]
分子分母同号，且在此之后 $y'<0$，所以 Mi 先升后降并最终趋于0。下图选择 $\mu=2\sigma$，横轴为无量纲时间 $s=\sigma t$，只是符合原条件的一个参数例子。
\begin{center}\begin{tikzpicture}\begin{axis}[width=.75\linewidth,height=4.7cm,axis lines=left,xlabel={$s$},ylabel={相对量},domain=0:6,samples=100,legend pos=north east]
\addplot[blue]{exp(-x)};\addlegendentry{St}
\addplot[red]{(exp(-x)-exp(-2*x))/2};\addlegendentry{Mi}
\addplot[black]{1-1.5*exp(-x)+.5*exp(-2*x)};\addlegendentry{En}
\end{axis}\end{tikzpicture}\end{center}
若另行考虑相等半衰期的极限 $\mu\to\sigma$，积分式连续给出 $y=(\sigma t/2)e^{-\sigma t}$、$z=1-(1+\sigma t/2)e^{-\sigma t}$，峰值时间为 $1/\sigma$；这不是本题给定条件，而是可用于检验一般公式的补充。
(e) 方程组线性且齐次，把初始总量加倍便使 $x,y,z$ 全部加倍。
(f) 代入得 $q=(t+2)e^t$。在 $t\ne0$ 的区间，齐次方程 $th'+2h=0$ 给 $h=C/t^2$，故 $x=e^t+C/t^2$。若要求在包含0的区间上连续可微，则 $C=0$，且 $t=0$ 处的原方程也要求 $x(0)=1$。
\end{solution}
\begin{exercise}\label{cw:ps2-II-5}
原题 II.5。(a) 制作三列表格，分别给直角式、模与辐角以及复平面上的点：(i) $1-i$；(ii) 模为2、辐角为 $\pi/6$；(iii) 虚部为负的 $i$ 的平方根；(iv) 辐角在 $(0,\pi/2)$ 的1的六次根；(v) $((1+i)/\sqrt2)^{-13}$。(b) 求所有根：(i) $z^4+4=0$；(ii) $z^2+2z+2=0$。
\end{exercise}
\begin{solution}
解答依据（编者译审或补解，AI 辅助）：官方 PS2 解答第2页。辐角均按模 $2\pi$ 理解；主辐角选 $(-\pi,\pi]$。
\begin{center}\begin{tabular}{ccl}序号&直角式&模与主辐角\\\hline
(i)&$1-i$&$\sqrt2,-\pi/4$\\
(ii)&$\sqrt3+i$&$2,\pi/6$\\
(iii)&$(-1-i)/\sqrt2$&$1,-3\pi/4$\\
(iv)&$(1+\sqrt3 i)/2$&$1,\pi/3$\\
(v)&$(-1+i)/\sqrt2$&$1,3\pi/4$\\
\end{tabular}\end{center}
(iii) 因 $i=e^{i\pi/2}$，两个平方根为 $\pm e^{i\pi/4}$，其中负号给负虚部。(iv) 六次根辐角为 $k\pi/3$，指定范围只允许 $k=1$。(v) 底数是 $e^{i\pi/4}$，负13次方为 $e^{-13i\pi/4}=e^{3i\pi/4}$。相应点如下；它们的坐标就是表格直角式的实、虚部。
\begin{center}\begin{tikzpicture}[scale=1.3]
\draw[->](-1.2,0)--(2.2,0)node[right]{实部};\draw[->](0,-1.25)--(0,1.5)node[above]{虚部};
\foreach \xx/\yy/\lab in {1/-1/i,1.732/1/ii,-.707/-.707/iii,.5/.866/iv,-.707/.707/v}{\fill(\xx,\yy)circle(1.5pt)node[above right]{(\lab)};}
\end{tikzpicture}\end{center}
(b)(i) $z^4=-4=4e^{i(\pi+2k\pi)}$，故 $z=\sqrt2e^{i(\pi/4+k\pi/2)}$，$k=0,1,2,3$，即 $1+i,-1+i,-1-i,1-i$。(ii) 完成平方得 $(z+1)^2=-1$，根为 $-1\pm i$。
\end{solution}
\begin{exercise}\label{cw:ps2-II-6}
原题 II.6。(a) 在 II.5(a) 的表中加一列指数式 $Ae^{i\theta}$。(b) 求所有满足 $e^{a+bi}=-2$ 的复数。(c) 由 $(e^{it})^4=e^{4it}$ 导出以 $\sin t,\cos t$ 的幂表示 $\cos4t$ 的公式。(d) 对下列 $f$，画实函数图、找 $w=a+bi$ 使 $\Re(e^{wt})=f(t)$，并画 $e^{wt}$ 的轨迹：(i) $\cos2\pi t$；(ii) $e^{-t}$；(iii) $e^{-t}\cos2\pi t$；(iv) 恒等于1。原题建议借助 Complex Exponential Mathlet。
\end{exercise}
\begin{solution}
解答依据（编者译审或补解，AI 辅助）：官方 PS2 解答第2页。(a) 五个指数式依次为 $\sqrt2e^{-i\pi/4}$、$2e^{i\pi/6}$、$e^{-3i\pi/4}$、$e^{i\pi/3}$、$e^{3i\pi/4}$。
(b) 取模得 $e^a=2$，故 $a=\ln2$；取单位复数部分得 $e^{ib}=-1$，故 $b=(2k+1)\pi$、$k\in\mathbb Z$，这给出全部解。
(c) 展开并取实部：
\[
 \cos4t=\Re(\cos t+i\sin t)^4
 =\cos^4t-6\cos^2t\sin^2t+\sin^4t.
\]
(d) 利用 $e^{(a+bi)t}=e^{at}(\cos bt+i\sin bt)$，可分别取 $w=2\pi i,-1,-1+2\pi i,0$。轨迹分别是单位圆、正实轴、绕原点逆时针向内的螺线、单点1；这里“向内”指 $t$ 增大。若只取 $t\geq0$，第二项轨迹为 $(0,1]$，第三项只取螺线从1向内的部分。下图画 $0\leq t\leq2$ 的实函数与相应轨迹，恒值1与实轴上单点重合。
\begin{center}\begin{tikzpicture}
\begin{axis}[name=realplot,width=.46\linewidth,height=4.3cm,axis lines=middle,xlabel={$t$},domain=0:2,samples=140,ymin=-1.15,ymax=1.2]
\addplot[blue]{cos(deg(2*pi*x))};\addplot[red]{exp(-x)};\addplot[black]{exp(-x)*cos(deg(2*pi*x))};\addplot[orange]{1};
\end{axis}
\begin{axis}[at={(realplot.east)},anchor=west,xshift=1cm,width=.42\linewidth,height=4.3cm,axis equal image,axis lines=middle,xlabel={实部},ylabel={虚部},xlabel style={at={(axis description cs:1.03,.5)},anchor=west},ylabel style={at={(axis description cs:.5,1.03)},anchor=south},xmin=-1.15,xmax=1.15,ymin=-1.15,ymax=1.15,samples=150]
\addplot[blue,domain=0:360]({cos(x)},{sin(x)});
\addplot[red,domain=0:2]({exp(-x)},0);
\addplot[black,domain=0:2]({exp(-x)*cos(deg(2*pi*x))},{exp(-x)*sin(deg(2*pi*x))});\addplot[orange,only marks]coordinates{(1,0)};
\end{axis}\end{tikzpicture}\end{center}
\end{solution}
\begin{exercise}\label{cw:ps2-II-7}
原题 II.7。(a) 把 $\Re(e^{3it}/(\sqrt3+i))$ 写成 $a\cos3t+b\sin3t$ 和 $A\cos(3t-\phi)$，再通过分母的极式重算。原题第二条路线中出现了不一致的 $2+2i$ 与 $e^{2it}$；请辨明对应表达式。(b) 为 $z'+3z=e^{2it}$ 找形如 $we^{2it}$ 的特解并写通解。(c) 用(b)求 $x'+3x=\cos2t$ 的特解和通解。
\end{exercise}
\begin{solution}
解答依据（编者译审或补解，AI 辅助）：官方 PS2 解答第2页，修正题面(a)第二路线误植。(a) 有
\[
 \frac1{\sqrt3+i}=\frac{\sqrt3-i}{4},\qquad
 \Re\frac{e^{3it}}{\sqrt3+i}=\frac{\sqrt3}4\cos3t+\frac14\sin3t
 =\frac12\cos(3t-\pi/6).
\]
另一方面 $\sqrt3+i=2e^{i\pi/6}$，于是商为 $\frac12e^{i(3t-\pi/6)}$，取实部得到同一结果。题面的 $2+2i$ 是 $2\sqrt2e^{i\pi/4}$，若真的计算 $e^{2it}/(2+2i)$，实部为 $\frac1{2\sqrt2}\cos(2t-\pi/4)$，不会等于本题前一个表达式，因此不能把两式混用。
(b) 代入得 $(3+2i)w=1$，所以 $z=e^{2it}/(3+2i)+Ce^{-3t}$，$C\in\C$。(c) 取实部得
\[
 x=\frac{3\cos2t+2\sin2t}{13}+ce^{-3t},\qquad c\in\R.
\]
求导代回即可核验右端确为 $\cos2t$。
\end{solution}

\originalchapter{作业三：平衡与阻尼}\label{cw:ps3}
本章合并官方 PS3 的前后两半。原讲次10是考试，因此题号从9跳到11；不重编号。
\originalsection{第一部分：指定练习}
\begin{exercise}\label{cw:ps3-I-8}原题 I.8：Notes 1E-1，见\hyperref[cw:bank-1E-1]{题1E-1}。\end{exercise}
\begin{solution}题面与完整解答见链接。\end{solution}
\begin{exercise}\label{cw:ps3-I-9}
原题 I.9。考虑 $x'+2x=1$。(a) 分别用(i)分离变量，(ii)积分因子，(iii)把1写为 $e^{0t}$、试常数特解后加齐次解，求通解。(b) 画相线和包括平衡解在内的若干解，判断稳定性。(c) 对 $x(0)=0$，用三个等长欧拉步估计 $x(1)$。
\end{exercise}
\begin{solution}
解答依据（编者译审或补解，AI 辅助）：PS3前半第2页部分答案，补全三条推导。(a)(i) 非平衡时 $dx/(1-2x)=dt$，积分得 $-\frac12\ln|1-2x|=t+C$，恢复被除去的平衡解后，$x=\frac12+ce^{-2t}$。(ii) 乘 $e^{2t}$ 后有 $(e^{2t}x)'=e^{2t}$，积分得到同一通解。(iii) 常数特解 $A$ 满足 $2A=1$，加上齐次解 $ce^{-2t}$ 即得结果。(b) 在 $x<1/2$ 时 $x'>0$，在 $x>1/2$ 时 $x'<0$，相线箭头都指向 $1/2$；任意解与 $1/2$ 的差为 $ce^{-2t}\to0$，故平衡是渐近稳定的。
\begin{center}\begin{tikzpicture}[scale=1.1]
\draw[->](-2,0)--(2,0)node[right]{$x$};\fill(0,0)circle(2pt)node[below]{$1/2$};\draw[->,thick](-1.6,.15)--(-.5,.15);\draw[->,thick](1.6,.15)--(.5,.15);
\end{tikzpicture}\end{center}
下图另取 $C=-1,-1/2,0,1/2,1$ 描出 $x(t)=1/2+Ce^{-2t}$，其中水平虚线为平衡解，所有曲线沿时间正向趋于它。
\begin{center}\begin{tikzpicture}\begin{axis}[width=.72\linewidth,height=4.3cm,axis lines=middle,xlabel={$t$},ylabel={$x$},domain=0:3,samples=100,ymin=-.6,ymax=1.6]
\addplot[blue]{.5-exp(-2*x)};\addplot[red]{.5-exp(-2*x)/2};
\addplot[black,dashed]{.5};\addplot[red]{.5+exp(-2*x)/2};\addplot[blue]{.5+exp(-2*x)};
\end{axis}\end{tikzpicture}\end{center}
(c) $h=1/3$，$x_{k+1}=x_k+h(1-2x_k)=x_k/3+1/3$，依次得 $x_1=1/3,x_2=4/9,x_3=13/27$。真值是 $(1-e^{-2})/2$，欧拉估计 $13/27$ 为所求。
\end{solution}
\begin{exercise}\label{cw:ps3-I-11}原题 I.11：Notes 2C-1(a)、(b)，见\hyperref[cw:bank-2C-1]{题2C-1}；EP 第2.1节第43、44、47题，第2.3节第1、2、7、15题。\end{exercise}
\begin{solution}公开 Notes 指定小问复用链接；EP 七道题只列定位，未取得题面，未翻译、未解答。\end{solution}
\begin{exercise}\label{cw:ps3-I-12}原题 I.12：Notes 2C-1(c)、(d)、(e)以及2G-1，见\hyperref[cw:bank-2C-1]{题2C-1}、\hyperref[cw:bank-2G-1]{题2G-1}；EP 第2.3节第5、21、23题。\end{exercise}
\begin{solution}公开题按链接复用；EP 三道题仅列定位，未翻译、未解答。\end{solution}
\originalsection{第二部分：题面与解答}
\begin{exercise}\label{cw:ps3-II-8}
原题 II.8。没有捕猎时，保护区羚羊服从稳定种群为1千只的逻辑斯谛模型；政府希望研究恒定年捕猎率 $a$ 的影响。原题用 Phase Lines Mathlet 的首项
\[
 y'=(1-y)y-a
\]
和相线、分岔图研究此问题。
(a) 说明此模型的假设。(b) 对所有 $a$ 求平衡点数量、位置与稳定性。(c) 若每年允许捕杀187.5只，稳定种群是多少？低于什么临界种群会崩溃？(d) 对此 $a$ 描出五种解行为；相差时间平移的解视为同型。说明向前和向后时间的行为，并与相线各部分对应。(e) 给出分岔图中平衡点曲线的方程。
\end{exercise}
\begin{solution}
解答依据（编者译审或补解，AI 辅助）：官方 PS3 解答第1页，修正其(e)符号笔误，补充五类解的端点行为。
(a) 一般逻辑斯谛增长为 $y'=k_0(1-y/p)y$。这里 $p=1$ ko，另取 $k_0=1$ 年$^{-1}$，并假设捕猎率不随种群大小变化，因而减去常数 $a$（ko/年）。这一模型只在种群非负时有生物意义，抵达0以后应停止而不继续把负值解释为种群。
(b) 平衡条件 $y^2-y+a=0$ 给
\[
 y_\pm=\frac12\pm\sqrt{\frac14-a}.
\]
$a<1/4$ 时有两个平衡点。由于 $F(y)=y-y^2-a=-(y-y_-)(y-y_+)$，外侧 $F<0$、两根之间 $F>0$，故低根不稳定、高根渐近稳定。$a=1/4$ 时 $F=-(y-1/2)^2$，唯一平衡从上方吸引、从下方排斥，是半稳定。$a>1/4$ 时没有平衡，所有解均下降。若讨论捕猎，只允许 $a\geq0$；代数上的 $a<0$ 表示恒定引入种群。
(c) $a=187.5/1000=3/16$，所以 $y_-=1/4,y_+=3/4$，稳定值750只，崩溃阈值250只。
(d) 相线分为 $y<1/4$、$y=1/4$、$1/4<y<3/4$、$y=3/4$、$y>3/4$ 五部分。两个平衡分别给恒定解。其余三部分用分离变量统一描述：
\[
 \frac{y-3/4}{y-1/4}=Ke^{-t/2},\qquad
 y(t)=\frac{3/4-(K/4)e^{-t/2}}{1-Ke^{-t/2}}.
\]
取 $K<0$ 时，解在两根之间，向后趋于 $1/4$，向前趋于 $3/4$。取 $K>0$ 时，分母零点 $T=2\ln K$ 分开两个最大解区间：在 $t>T$，解从 $+\infty$ 下降到 $3/4$；在 $t<T$，解从 $1/4$ 的下方下降，$t\uparrow T$ 时趋于 $-\infty$，且在此前已抵达种群0。如下图分别选 $K=1,-1$，画数学模型三类非恒定解和两条平衡解。
\begin{center}\begin{tikzpicture}\begin{axis}[width=.76\linewidth,height=5.2cm,axis lines=middle,xlabel={$t$},ylabel={$y$},xmin=-5,xmax=5,ymin=-.5,ymax=1.5,samples=100]
\addplot[blue,domain=.45:5]{(.75-.25*exp(-x/2))/(1-exp(-x/2))};
\addplot[red,domain=-5:-.45]{(.75-.25*exp(-x/2))/(1-exp(-x/2))};
\addplot[black,domain=-5:5]{(.75+.25*exp(-x/2))/(1+exp(-x/2))};
\addplot[gray,dashed,domain=-5:5]{.25};\addplot[gray,dashed,domain=-5:5]{.75};
\end{axis}\end{tikzpicture}\end{center}
(e) 分岔曲线为 $a=y-y^2$，等价于 $y^2-y+a=0$。官方答案把最后一项误写为 $-a$；代回原方程可立即发现此误。两分支在 $(a,y)=(1/4,1/2)$ 合并。
\end{solution}
\begin{exercise}\label{cw:ps3-II-9}
原题 II.9。仍取 $y'=(1-y)y-3/16$，令稳定平衡为 $y_0$，$u=y-y_0$ 表示种群相对平衡的超额。(a) 写出 $u$ 的方程，检验自治性与零平衡。(b) 忽略 $u^2$ 等高阶项，求线性化方程及通解。(c) 写出近平衡近似；若 $y(10)-y_0=b$，估计 $y(11),y(12)$。(d) 若线性方程 $x'+p(t)x=q(t)$ 是自治方程，$p,q$ 应满足什么条件？
\end{exercise}
\begin{solution}
解答依据（编者译审或补解，AI 辅助）：官方 PS3 解答第1页，补上(c)的两个具体预测和(d)的理由。$y_0=3/4$，代入 $y=u+3/4$ 得
\[
 u'=(1/4-u)(u+3/4)-3/16=-\tfrac12u-u^2.
\]
右端无显含时间，且在 $u=0$ 时为0。(b) 线性化为 $u'=-u/2$，故 $u=Ce^{-t/2}$。(c) 当 $|u|\ll1$ 时二次项相对线性项很小，得到
\[
 y(t)\simeq\tfrac34+be^{-(t-10)/2},\quad
 y(11)\simeq\tfrac34+be^{-1/2},\quad y(12)\simeq\tfrac34+be^{-1}.
\]
这些是近平衡预测，不是原非线性方程的精确等式。精确超额也可写为
\[
 u(t)=\frac{be^{-(t-10)/2}}{1+2b(1-e^{-(t-10)/2})},
\]
在分母非零的区间成立；它显示被忽略的修正从 $b^2$ 开始。(d) 自治性要求 $q(t)-p(t)x$ 对每个 $x$ 都不依赖 $t$。取 $x=0$ 得 $q$ 常数，再取 $x=1$ 得 $p$ 也是常数。
\end{solution}
\begin{exercise}\label{cw:ps3-II-11}
原题 II.11。酒馆门可穿过门框，阻尼应使门缓慢回到闭合位置；若向门框推得足够快，它可能穿过后从另一侧回到闭合位置。用 Damped Vibration Mathlet 研究
\[
 mx''+bx'+kx=0,
\]
初值框的横坐标是 $x'(0)$、纵坐标是 $x(0)$。固定 $m=0.50,b=1.5,k=0.625=5/8$；Mathlet 可只能显示 $k=0.62$，计算仍用精确值。(a) 求特征多项式与根。(b) 用 $x(0)=x_0,x'(0)=v_0$ 写通解。(c) 哪些初值给出单一指数解？(d) 固定 $x_0$（示例 $x_0=0.25$），求使解在某个正时间等于0的 $v_0$ 的精确范围。
\end{exercise}
\begin{solution}
解答依据（编者译审或补解，AI 辅助）：官方 PS3 解答第2页，补充 $x_0=0$ 的边界。(a) 特征多项式
\[
 p(r)=\tfrac12r^2+\tfrac32r+\tfrac58
 =\tfrac12(r+\tfrac12)(r+\tfrac52),
\]
根为 $-1/2,-5/2$。(b) 设 $x=c_1e^{-t/2}+c_2e^{-5t/2}$，则 $c_1+c_2=x_0$，$-c_1/2-5c_2/2=v_0$。解这两个线性方程得
\[
 c_1=\tfrac14(5x_0+2v_0),\quad c_2=-\tfrac14(x_0+2v_0).
\]
(c) $c_1=0$ 或 $c_2=0$，即 $v_0=-5x_0/2$ 或 $v_0=-x_0/2$；零解位于两直线交点。
(d) 非零解有零点时，$e^{2t}=-c_2/c_1$。正时间零点要求 $-c_2/c_1>1$。若 $x_0>0$，这等价于 $c_1<0$，即 $v_0<-5x_0/2$；若 $x_0<0$，等价于 $v_0>-5x_0/2$。所以 $x_0=1/4$ 时须 $v_0<-5/8$。等号时是单一快指数解，只渐近到0，不在有限时间穿过。若 $x_0=0$ 且 $v_0\ne0$，$x=(v_0/2)(e^{-t/2}-e^{-5t/2})$，在 $t>0$ 没有零点；$x_0=v_0=0$ 时是恒零解。门的位置与速度可直接从公式重绘，无需把交互读数当作计算依据。
\end{solution}
\begin{exercise}\label{cw:ps3-II-12}
原题 II.12。继续使用同一模型，保持 $m=0.50,k=0.625$。(a) 降低 $b$，根在何值重合？代数验证并给此时通解。(b) 再取 $b=0.25$，任选非零初值，求根与通解。(c) 对 $x(0)=0,x'(0)=1$，求前五个零点的相邻间距和伪周期，判断振动是否随衰减而加快或减慢；再对不同初值说明相邻零点间距。
\end{exercise}
\begin{solution}
解答依据（编者译审或补解，AI 辅助）：官方 PS3 解答第2页，提供理论零点以替代未经运行的测量。(a) 归一化多项式为 $r^2+2br+5/4$，判别式为 $4b^2-5$。非负阻尼的重根条件是 $b=\sqrt5/2$，根为 $-\sqrt5/2$，通解 $(A+Bt)e^{-\sqrt5t/2}$。
(b) 当 $b=1/4$，多项式可写为 $(r+1/4)^2+19/16$，所以
\[
 r=-\tfrac14\pm i\tfrac{\sqrt{19}}4,\qquad
 x=e^{-t/4}\left(A\cos\tfrac{\sqrt{19}}4t+B\sin\tfrac{\sqrt{19}}4t\right).
\]
(c) 给定初值使 $A=0,B=4/\sqrt{19}$，零点为 $t_j=4j\pi/\sqrt{19}$，$j=0,1,2,\ldots$。含 $t=0$ 的前五个为
\[
 0,\quad\frac{4\pi}{\sqrt{19}},\quad\frac{8\pi}{\sqrt{19}},\quad\frac{12\pi}{\sqrt{19}},\quad\frac{16\pi}{\sqrt{19}}.
\]
若只数正零点，则从 $j=1$ 开始。每个相邻间距均为 $4\pi/\sqrt{19}\simeq2.8829$，伪周期为其两倍 $8\pi/\sqrt{19}\simeq5.7658$。非零一般解写成 $Re^{-t/4}\cos(\sqrt{19}t/4-\phi)$ 后，零点只整体平移，相邻间距保持相同。振幅衰减，振动的角频率和伪周期保持不变。
\begin{center}\begin{tikzpicture}\begin{axis}[width=.75\linewidth,height=4.6cm,axis lines=middle,xlabel={$t$},ylabel={$x$},domain=0:13,samples=180]
\addplot[blue]{4/sqrt(19)*exp(-x/4)*sin(deg(sqrt(19)*x/4))};
\addplot[gray,dashed]{4/sqrt(19)*exp(-x/4)};\addplot[gray,dashed]{-4/sqrt(19)*exp(-x/4)};
\end{axis}\end{tikzpicture}\end{center}
\end{solution}

\originalchapter{作业四：受迫响应与共振}\label{cw:ps4}
\originalsection{第一部分：指定练习}
\begin{exercise}\label{cw:ps4-I-13}原题 I.13：Notes 2F-6(b)、(c)、(d)，见\hyperref[cw:bank-2F-6]{题2F-6}。\end{exercise}
\begin{solution}仅复用指定三个小问的公开题面与解答。\end{solution}
\begin{exercise}\label{cw:ps4-I-14}原题 I.14：求 $x'''-x=e^{2t}$ 的通解。\end{exercise}
\begin{solution}
解答依据（编者译审或补解，AI 辅助）：PS4题面第3页答案。特征多项式 $r^3-1=(r-1)(r^2+r+1)$，根为 $1,-1/2\pm i\sqrt3/2$。试特解 $Be^{2t}$ 得 $7B=1$，因此
\[
 x=\tfrac17e^{2t}+c_1e^t+e^{-t/2}
 \left(c_2\cos\tfrac{\sqrt3}2t+c_3\sin\tfrac{\sqrt3}2t\right).
\]
\end{solution}
\begin{exercise}\label{cw:ps4-I-15}原题 I.15：Notes 2C-8(c)、(d)，见\hyperref[cw:bank-2C-8]{题2C-8}；EP 第2.5节第2、8、11、14题。\end{exercise}
\begin{solution}公开题按链接复用；EP 四道题仅列定位，题面未取得，未翻译、未解答。\end{solution}
\begin{exercise}\label{cw:ps4-I-16}原题 I.16：EP 第2.6节第17题。\end{exercise}
\begin{solution}仅列定位，题面未取得，未翻译、未解答。\end{solution}
\originalsection{第二部分：题面与解答}
\begin{exercise}\label{cw:ps4-II-13}
原题 II.13。(a) 求 $x'+2x=e^{3t}\cos4t$ 的实通解，建议用复数与指数响应公式。(b) 求 $x''+x'+2x=\cos t$ 的实通解，把稳态特解分别写成 $a\cos t+b\sin t$ 和 $A\cos(t-\phi)$。
原题随后使用 Forced Damped Vibrations Mathlet，力直接作用于质量，固定 $m=b=1,k=2$。(c) 在输入幅值 $A=0$、初值 $x(0)=0,x'(0)=1$ 下求伪周期。(d) 取输入 $A=1,\omega=1$，求稳态幅值、$x_p(0),x_p'(0)$，与图比较。(e) 对静止初值 $x(0)=x'(0)=0$，求瞬态项 $x_h$，使 $x=x_p+x_h$。
\end{exercise}
\begin{solution}
解答依据（编者译审或补解，AI 辅助）：官方 PS4 解答第1页。(a) 复特解 $z_p=e^{(3+4i)t}/(5+4i)$，分母有理化给
\[
 x=\frac{e^{3t}}{41}(5\cos4t+4\sin4t)+ce^{-2t}.
\]
(b) 特征多项式 $p(r)=r^2+r+2$，$p(i)=1+i$，故
\[
 x_p=\Re\frac{e^{it}}{1+i}=\tfrac12(\cos t+\sin t)
 =\frac1{\sqrt2}\cos(t-\pi/4).
\]
齐次根为 $-1/2\pm i\sqrt7/2$，实通解为
\[
 x=x_p+e^{-t/2}(C\cos\tfrac{\sqrt7}2t+D\sin\tfrac{\sqrt7}2t).
\]
(c) 无输入且给定初值时，$x=(2/\sqrt7)e^{-t/2}\sin(\sqrt7t/2)$，因此伪周期 $4\pi/\sqrt7\simeq4.7496$；这是理论值，不是交互测量值。(d) 幅值 $1/\sqrt2\simeq0.7071$，初值 $x_p(0)=x_p'(0)=1/2$。
(e) 需使 $x_h(0)=x_h'(0)=-1/2$。从 $C=-1/2$ 和 $-C/2+\sqrt7D/2=-1/2$ 得 $D=-3/(2\sqrt7)$，所以
\[
 x_h=e^{-t/2}\left(-\tfrac12\cos\tfrac{\sqrt7}2t
 -\frac3{2\sqrt7}\sin\tfrac{\sqrt7}2t\right).
\]
两项相加在 $t=0$ 的位移和速度均为0。图按所得公式复现三条曲线：
\begin{center}\begin{tikzpicture}\begin{axis}[width=.76\linewidth,height=5cm,axis lines=middle,xlabel={$t$},domain=0:12,samples=200,legend pos=north east]
\addplot[green!50!black]{(cos(deg(x))+sin(deg(x)))/2};\addlegendentry{稳态}
\addplot[red]{exp(-x/2)*(-cos(deg(sqrt(7)*x/2))/2-3*sin(deg(sqrt(7)*x/2))/(2*sqrt(7)))};\addlegendentry{瞬态}
\addplot[blue]{(cos(deg(x))+sin(deg(x)))/2+exp(-x/2)*(-cos(deg(sqrt(7)*x/2))/2-3*sin(deg(sqrt(7)*x/2))/(2*sqrt(7)))};\addlegendentry{总解}
\end{axis}\end{tikzpicture}\end{center}
\end{solution}
\begin{exercise}\label{cw:ps4-II-14}
原题 II.14。Amplitude and Phase: Second Order I Mathlet 研究经弹簧驱动的单位质量系统。输入为弹簧上端的位置 $y(t)=\cos\omega t$，输出为质量位置的稳态 $x(t)$，输入幅值为1。(a) 取 $b=0.5,k=4,\omega=2$，核算增益和时间滞后。(b) 求复增益、随 $\omega$ 变化的增益，以及相位滞后 $\phi$ 的正切。
\end{exercise}
\begin{solution}
解答依据（编者译审或补解，AI 辅助）：官方 PS4 解答第1页。弹簧力为 $k(y-x)$，阻尼力为 $-bx'$，故 $x''+bx'+kx=ky$，输入经弹簧而使右端带因子 $k$。
(a) $p(2i)=4-4+i=i$，于是 $x_p=\Re(4e^{2it}/i)=4\sin2t$，增益4，相位滞后 $\phi=\pi/2$，时间滞后为 $\phi/\omega=\pi/4\simeq0.7854$。
(b) 复输入 $e^{i\omega t}$ 对应稳态 $H(\omega)e^{i\omega t}$，其中
\[
 H(\omega)=\frac4{4-\omega^2+i\omega/2},\quad
 g(\omega)=\frac4{\sqrt{(4-\omega^2)^2+\omega^2/4}},\quad
 \tan\phi=\frac{\omega/2}{4-\omega^2}.
\]
$\omega>0$ 时应取 $\phi=\operatorname{Arg}(4-\omega^2+i\omega/2)\in(0,\pi)$；仅靠正切不能确定象限，$\omega=2$ 时正切无定义但相位仍为 $\pi/2$。
\end{solution}
\begin{exercise}\label{cw:ps4-II-15}
原题 II.15。求下列方程的实通解：(a) $x'+x=e^{-t}$，使用共振指数响应公式；(b) $x'''-x'=t^2+1$；(c) $2x''+x'=(t^2+1)e^{2t}$。
\end{exercise}
\begin{solution}
解答依据（编者译审或补解，AI 辅助）：(a)(b)官方 PS4 解答第1页；(c)官方文件未给解答，独立补算。
(a) $r+1$ 在 $r=-1$ 为零，因此试 $te^{-t}$。其导数加自身为 $e^{-t}$，通解 $x=(t+C)e^{-t}$。
(b) 令 $u=x'$，则 $u''-u=t^2+1$。试 $u=at^2+bt+c$，比较系数 $-a=1,-b=0,2a-c=1$，得 $u_p=-t^2-3$。积分得 $x_p=-t^3/3-3t$；齐次多项式 $r(r-1)(r+1)$ 给
\[
 x=-t^3/3-3t+C_1+C_2e^t+C_3e^{-t}.
\]
(c) 为避免反复求带指数的导数，令 $x=e^{2t}v$。则
\[
 2x''+x'=e^{2t}(2v''+9v'+10v).
\]
试 $v=At^2+Bt+C$，比较 $10A=1$、$18A+10B=0$、$4A+9B+10C=1$，得 $A=1/10$、$B=-9/50$、$C=111/500$。齐次根为0、$-1/2$，所以
\[
 x=e^{2t}\left(\frac{t^2}{10}-\frac{9t}{50}+\frac{111}{500}\right)
 +C_1+C_2e^{-t/2}.
\]
这里常数项经 $4A+9B+10C$ 核验为1。
\end{solution}
\begin{exercise}\label{cw:ps4-II-16}
原题 II.16。继续用 II.14 的经弹簧驱动系统，前三问固定 $k=4,b=0.5$。(a) 画增益随输入角频率的图，求使增益最大的实际共振频率 $\omega_r$ 和最大增益，判断是否恰为2。(b) 求相位滞后45度时的正频率。(c) 画 $\omega$ 从0到无穷时的复增益 Nyquist 轨迹，并标出相位接近 $\pi/4$ 的点。(d) 固定 $\omega=2,k=4$、改变 $b$，复增益在图中如何移动？对一般 $k$，调整 $\omega$ 使相位 $\pi/2$ 后再改变 $b$，相位为何保持？给出一般 $H$ 并解释。
\end{exercise}
\begin{solution}
解答依据（编者译审或补解，AI 辅助）：官方 PS4 答案文件没有本题解答，以下独立由 II.14 的公式求解。
(a) 最大化 $g$ 等价于最小化
\[
 D(\omega)=(4-\omega^2)^2+\omega^2/4,\qquad
 D'(\omega)=\omega(4\omega^2-31/2).
\]
正临界点 $\omega_r=\sqrt{31/8}\simeq1.9685$，在它之前 $D$ 下降、之后上升，所以给出唯一最大增益。此时 $4-\omega_r^2=1/8$，$D=63/64$，从而 $g_{\max}=32/\sqrt{63}\simeq4.0316$。它略高于 $g(2)=4$。
(b) 相位 $\pi/4$ 要求实部、虚部相等且为正，$4-\omega^2=\omega/2$，正解
\[
 \omega_{45}=\frac{\sqrt{65}-1}{4}\simeq1.7656.
\]
(c) 有
\[
 \Re H=\frac{4(4-\omega^2)}{D(\omega)},\qquad
 \Im H=-\frac{2\omega}{D(\omega)}.
\]
轨迹从 $H(0)=1$ 出发，经第四象限到 $H(2)=-4i$，随后进入第三象限，最终从下方趋于0。在 $\omega_{45}$ 处，$H=(4/\omega_{45})(1-i)$，下图标出该点，并同时给出增益图。
\begin{center}\begin{tikzpicture}
\begin{axis}[name=gainplot,width=.45\linewidth,height=4.7cm,axis lines=left,xlabel={$\omega$},ylabel={增益},domain=0:4,samples=220]
\addplot[blue]{4/sqrt((4-x*x)^2+x*x/4)};\addplot[only marks,red]coordinates{(1.968502,4.031621)};
\end{axis}
\begin{axis}[at={(gainplot.east)},anchor=west,xshift=1.1cm,width=.42\linewidth,height=4.7cm,axis equal image,axis lines=middle,xlabel={$\Re H$},ylabel={$\Im H$},xmin=-2.4,xmax=3.1,ymin=-4.3,ymax=.3,samples=250]
\addplot[blue,domain=0:10]({4*(4-x*x)/((4-x*x)^2+x*x/4)},{-2*x/((4-x*x)^2+x*x/4)});
\addplot[red,only marks]coordinates{(2.265564,-2.265564)};
\draw[red,->](axis cs:0,0)--(axis cs:2.265564,-2.265564);
\end{axis}\end{tikzpicture}\end{center}
(d) 一般的单位质量系统有
\[
 H(\omega)=\frac{k}{k-\omega^2+ib\omega}.
\]
对 $b>0,k>0,\omega>0$，$\phi=\pi/2$ 意味着 $H$ 是负纯虚数，等价于 $k-\omega^2=0$。所以 $\omega=\sqrt k$ 时无论正阻尼取何值，相位都保持 $\pi/2$，且 $H=-i\sqrt k/b$ 沿负虚轴移动。特别地 $k=4,\omega=2$ 时 $H=-2i/b$。$b=0$ 是无阻尼共振，稳态公式分母为零，不能纳入此论断。
\end{solution}

\originalchapter{作业五：电路与傅里叶级数}\label{cw:ps5}
本章合并官方 PS5 的两半；第19讲为考试，第一部分第18项是“无新题”，这两项保留为安排说明而不另造题目。
\originalsection{第一部分：指定练习}
\begin{exercise}\label{cw:ps5-I-17}
原题 I.17：EP 第2.5节第8题；另外，求 $x''+x=t^2+\cos(2t-1)$ 的一个解。
\end{exercise}
\begin{solution}
EP 题仅列定位，未翻译、未解答；它曾出现在 PS4 I.15。对于显式给出的方程，依据 PS5 后半第1页官方更正独立核验。设多项式特解 $At^2+Bt+C$，则 $x''+x=At^2+Bt+(2A+C)$，比较得 $A=1,B=0,C=-2$。对余弦项，二阶导加自身使系数乘 $1-2^2=-3$，故一个解为
\[
 x_p=t^2-2-\tfrac13\cos(2t-1).
\]
前半第2页曾把常数写成 $-1/2$，后半明确更正为 $-2$。直接代回，$(t^2-2)''+(t^2-2)=t^2$，确认更正；若需要通解，再加 $C_1\cos t+C_2\sin t$。
\end{solution}
原题 I.18：无新题。I.19 与 II.19：课堂考试，不是额外作业题。
\begin{exercise}\label{cw:ps5-I-20}
原题 I.20：Notes 7A-1、7A-2(a)，见\hyperref[cw:bank-7A-1]{题7A-1}、\hyperref[cw:bank-7A-2]{题7A-2(a)}。
\end{exercise}
\begin{solution}中文题面与解答复用上述公开题。\end{solution}
\begin{exercise}\label{cw:ps5-I-21}
原题 I.21：(a) 再做 Notes 7A-2(a)，这次先以方波 $\operatorname{sq}(t)$ 表示该函数；(b) 做 Notes 7A-2(b)，通过积分其导数的傅里叶级数求解。两题见\hyperref[cw:bank-7A-2]{题7A-2}。
\end{exercise}
\begin{solution}
解答依据（编者补解，AI 辅助）：公开 Notes 7A-2 的题面，按本次作业另行指定的方法重算。
(a) 链接中的函数在 $-\pi<t\leq0$ 为0、在 $0<t\leq\pi$ 为1，随后按 $2\pi$ 周期延拓。在非跳点处它等于 $(1+\operatorname{sq}(t))/2$，因此
\[
 f(t)\sim\tfrac12+\frac2\pi\sum_{j=0}^\infty\frac{\sin((2j+1)t)}{2j+1}.
\]
原题在跳点所定的0或1与方波表示的跳点值 $1/2$ 不同；单点值不改变系数，级数在跳点收敛到左右极限的平均 $1/2$。不能把这个表示写成处处等式。
(b) 链接中的另一函数在 $[-\pi,\pi]$ 为 $|t|$，按周期延拓。除角点外，$f'=\operatorname{sq}(t)$。从方波级数积分得
\[
 f(t)=C-\frac4\pi\sum_{j=0}^\infty\frac{\cos((2j+1)t)}{(2j+1)^2}.
\]
其平均值由一周期内三角形面积除以周期得 $C=\pi/2$。积分的合法性可由方波部分和在 $L^2$ 中收敛、柯西--施瓦茨不等式控制原函数的积分误差说明（详见 II.21(f) 的同一论证）；所得 $n^{-2}$ 系数级数绝对一致收敛，在角点也收敛到连续函数 $|t|$ 的周期延拓值。
\end{solution}
\originalsection{第二部分：题面与解答}
\begin{exercise}\label{cw:ps5-II-17}
原题 II.17。串联 RLC 电路的电流满足
\[
 LI''+RI'+\frac1C I=V',\qquad V(t)=V_0\sin\omega t.
\]
采用 SI：$R$ 为欧姆、$L$ 为亨利、$C$ 为法拉、$V$ 为伏特、$I$ 为安培。原 Mathlet 的滑块用毫亨（$10^{-3}$ H）、微法（$10^{-6}$ F）和毫秒（$10^{-3}$ s）。这里只研究稳态电流，不研究各元件的电压降。
(a) 找出电流与电压同相的参数；特别地 $L=500$ mH、$\omega=200$ rad/s 时，$R,C,V_0$ 应如何选？(b) 推导同相的参数关系。(c) 取 $R=100\,\Omega,L=1000$ mH、$C=100\,\mu$F、$V_0=500$ V，求最大电流幅值所对应的频率、幅值与相位滞后。(d) 对一般参数计算并验证(c)的三项结论。
\end{exercise}
\begin{solution}
解答依据（编者译审或补解，AI 辅助）：官方 PS5 解答第1页；补全省漏的 $V_0$ 和阻尼项记号。令 $R,L,C,V_0>0$、$\omega>0$。为使输入相位清楚，以复电压 $V_0e^{i\omega t}$ 的虚部表示原输入。复稳态电流与复电压的比为
\[
 H_I(\omega)=\frac{i\omega}{1/C-L\omega^2+iR\omega}
 =\frac1{R+i(L\omega-1/(C\omega))}.
\]
电流为 $I_p=\Im(V_0H_Ie^{i\omega t})$。同相要求 $H_I$ 为正实数，所以必须且只须 $L\omega=1/(C\omega)$，即 $LC\omega^2=1$。
(a)(b) 在 $L=0.5$ H、$\omega=200$ rad/s 时，$C=1/(0.5\cdot200^2)=50\,\mu$F；任意正 $R,V_0$ 均同相，电流幅值是 $V_0/R$。例如取 $R=100\,\Omega,V_0=500$ V 就得到 $5\sin200t$ A。
(c)(d) 对一般参数，电流幅值与相位滞后为
\[
 A_I(\omega)=\frac{V_0}{\sqrt{R^2+(L\omega-1/(C\omega))^2}},\qquad
 \phi_I=\arctan\frac{L\omega-1/(C\omega)}R.
\]
分母至少为 $R$，且正频率下恰在 $\omega_r=1/\sqrt{LC}$ 取等，因此此处幅值最大为 $V_0/R$、滞后为0。低频 $\phi_I<0$ 表示电流超前，高频 $\phi_I>0$ 表示滞后。给定 $L=1$ H、$C=10^{-4}$ F，得到 $\omega_r=100$ rad/s、幅值5 A、电流电压同相。下图画幅频曲线和共振时的同相波形；横轴波形采用无量纲 $s=100t$，电压除以100以便和电流对照。
\begin{center}\begin{tikzpicture}
\begin{axis}[name=rlcamp,width=.45\linewidth,height=4.4cm,axis lines=left,xlabel={$\omega$ (rad/s)},ylabel={幅值 (A)},domain=1:250,samples=180]
\addplot[blue]{500/sqrt(10000+(x-10000/x)^2)};
\end{axis}
\begin{axis}[at={(rlcamp.east)},anchor=west,xshift=1cm,width=.43\linewidth,height=4.4cm,axis lines=middle,xlabel={$s$},domain=0:13,samples=180]
\addplot[blue,thick]{5*sin(deg(x))};\addplot[green!60!black,dashed]{5*sin(deg(x))};
\end{axis}\end{tikzpicture}\end{center}
\end{solution}
\begin{exercise}\label{cw:ps5-II-18}
原题 II.18。对 $mx''+bx'+kx=0$，自然角频率为 $\omega_n=\sqrt{k/m}$，临界阻尼满足 $b/m=2\omega_n$。定义阻尼比 $\zeta$ 使 $b/m=2\omega_n\zeta$。假设 $0<\zeta<1$，观察到经过 $n$ 个周期后指数包络减为原来的一半。用 $n$ 和 $\alpha=(\ln2)/(2\pi)\simeq0.1103178$ 求 $\zeta$；算 $n=10$，并与工程近似 $\alpha/n$ 比较。
\end{exercise}
\begin{solution}
解答依据（编者译审或补解，AI 辅助）：官方 PS5 解答第1页。欠阻尼解为
\[
 x=Ae^{-\zeta\omega_n t}\cos(\omega_dt-\phi),
 \qquad\omega_d=\omega_n\sqrt{1-\zeta^2}.
\]
$n$ 个伪周期历时 $2\pi n/\omega_d$，包络比为 $\exp[-2\pi n\zeta/\sqrt{1-\zeta^2}]$。令其等于 $1/2$，得到
\[
 \frac{\zeta}{\sqrt{1-\zeta^2}}=\frac\alpha n,
 \qquad\zeta=\frac{\alpha/n}{\sqrt{1+(\alpha/n)^2}}
 =\frac\alpha{\sqrt{n^2+\alpha^2}}.
\]
$n=10$ 时，精确值约 $0.0110311087$，近似 $\alpha/10\simeq0.0110317800$；近似略大，相对误差约 $0.006085\%$。包络峰值递减与同类局部峰值的递减比例一致：$x'$ 仍是同一频率的衰减正弦，极值等间隔，同号峰值间隔一个伪周期。
\end{solution}
\begin{exercise}\label{cw:ps5-II-20}
原题 II.20。Fourier Coefficients Mathlet 可用滑块设置正弦或余弦的有限和。选择目标[B]，先判断奇偶性，决定使用正弦还是余弦、奇数项还是偶数项，并解释。
(a) 用目测调整滑块，记录近似最优系数。(b) 目标实际为周期 $2\pi$ 的函数
\[
 f(t)=\begin{cases}\pi/4,&-\pi/2<t<\pi/2,\\-\pi/4,&\pi/2<t<3\pi/2,\end{cases}
 \qquad f(\pm\pi/2)=0,
\]
按周期延拓。用积分公式计算系数，与(a)比较。(c) 从任意滑块值开始，逐个调系数，使所示的均方根距离最小，记录最佳系数并与(b)比较；解释为什么最佳拟合就是傅里叶系数，且每个系数的优化互不影响。
\end{exercise}
\begin{solution}
解答依据（编者译审或补解，AI 辅助）：官方 PS5 解答第1--2页；补充(c)正交最小化证明。本文没有运行 Mathlet 或取得个人目测记录；(a)提供由公式可重现的比较值和有限和图，不能把这些数值冒称交互实验结果。
目标 $f$ 为偶函数，所以所有 $b_n=0$。又 $f(t+\pi)=-f(t)$，故平均值与偶数频率余弦系数为0：平移半周期后，偶数频率余弦不变，而函数改变符号，积分只能为0。因此只需奇数余弦。
(b) 用约定 $f\sim a_0/2+\sum_{n\geq1}(a_n\cos nt+b_n\sin nt)$。平均值0；由偶性得
\begin{align*}
 a_n&=\frac2\pi\left(\int_0^{\pi/2}\frac\pi4\cos nt\,dt
 -\int_{\pi/2}^{\pi}\frac\pi4\cos nt\,dt\right)\\
 &=\frac1{2n}\left(2\sin(n\pi/2)-\sin n\pi\right)
 =\frac{\sin(n\pi/2)}n.
\end{align*}
所以
\[
 f(t)\sim\cos t-\frac13\cos3t+\frac15\cos5t-\frac17\cos7t+\cdots.
\]
它在连续点收敛到 $f$，在跳点收敛到左右值平均0，与题目的跳点值一致。前五个非零系数是 $1,-1/3,1/5,-1/7,1/9$，用这些值可重绘(a)要求的有限和图。
(c) 给定有限和 $S=c_0+\sum_{n=1}^N(c_n\cos nt+d_n\sin nt)$，定义
\[
 E(S)^2=\frac1{2\pi}\int_{-\pi}^{\pi}|f-S|^2\,dt.
\]
三角函数正交：不同频率的乘积积分为0，每个正弦、余弦的平方积分为 $\pi$。展开平方并利用(b)中系数定义，得到
\[
 E(S)^2=\|f\|_{\rm rms}^2-c_0^*{}^2
 -\tfrac12\sum_{n=1}^N(a_n^2+b_n^2)
 +(c_0-c_0^*)^2+\tfrac12\sum_{n=1}^N[(c_n-a_n)^2+(d_n-b_n)^2],
\]
其中 $c_0^*=a_0/2=0$、$\|f\|_{\rm rms}^2=\pi^2/16$。每个可变系数只出现在自己的一项平方中，故唯一最优值是 $c_0=0,c_n=a_n,d_n=0$，调整一个不影响另一个的最佳值。比如保留 $n=1,3,5,7,9$ 时，最小均方根距离为
\[
 \left[\frac{\pi^2}{16}-\frac12\left(1+\frac19+\frac1{25}+\frac1{49}+\frac1{81}\right)\right]^{1/2}
 \simeq0.1578537.
\]
此数值是积分公式的结果，不是软件画面的读取。有限和在跳点附近有振铃，不影响平方积分的最佳拟合性质。
\begin{center}\begin{tikzpicture}\begin{axis}[width=.8\linewidth,height=5cm,axis lines=middle,xlabel={$t$},xmin=-3.1416,xmax=3.1416,ymin=-1.05,ymax=1.05,domain=-pi:pi,samples=220]
\addplot[green!50!black,thick]coordinates{(-3.14159,-.785398)(-1.570796,-.785398)(-1.570796,.785398)(1.570796,.785398)(1.570796,-.785398)(3.14159,-.785398)};
\addplot[blue,domain=-pi:pi]{cos(deg(x))-cos(deg(3*x))/3+cos(deg(5*x))/5-cos(deg(7*x))/7+cos(deg(9*x))/9};
\addplot[only marks,mark=*,green!50!black]coordinates{(-1.570796,0)(1.570796,0)};
\end{axis}\end{tikzpicture}\end{center}
图中竖直线段仅标示跳跃位置，并非函数在该横坐标具有多个值；跳点的实际函数值由实心点0标出。
\end{solution}
\begin{exercise}\label{cw:ps5-II-21}
原题 II.21。(a) 求 $2\sin(t-\pi/3)$ 的傅里叶级数。
定义奇的 $2\pi$ 周期方波 $\operatorname{sq}(t)$：在 $0<t<\pi$ 为1，在 $-\pi<t<0$ 为 $-1$，在整数倍 $\pi$ 处为0。
(b) 把同一方波视为 $4\pi$ 周期函数，用积分公式求其傅里叶级数，比较两种周期约定。
以下各项不得重新计算傅里叶系数积分，应从方波级数通过代数和微积分求级数、描图并给出最小周期：(c) $\operatorname{sq}(t-\pi/4)$；(d) $1+2\operatorname{sq}(2\pi t)$；(e) II.20 的 $f(t)$；(f) $2\pi$ 周期函数 $g$，在 $-\pi/2\leq t\leq\pi/2$ 时 $g=t$，在 $\pi/2\leq t\leq3\pi/2$ 时 $g=\pi-t$。提示：先联系 $g'$ 与 $f$。
\end{exercise}
\begin{solution}
解答依据（编者译审或补解，AI 辅助）：官方 PS5 解答第2页；(a)补回官方答案遗漏的题面系数2，(f)说明积分级数的依据。
(a) 角差公式直接给
\[
 2\sin(t-\pi/3)=-\sqrt3\cos t+\sin t.
\]
这已是傅里叶表达式，只含第一频率；由正交性表示唯一。官方答案写的是没有因子2的 $\sin(t-\pi/3)$，不能直接作为本题答案。
(b) 对周期 $2L=4\pi$，$L=2\pi$，基函数为 $\sin(nt/2),\cos(nt/2)$。奇性给 $a_n=0$，而
\begin{align*}
 b_n&=\frac1\pi\left(\int_0^\pi\sin(nt/2)\,dt-\int_\pi^{2\pi}\sin(nt/2)\,dt\right)\\
 &=\frac2{\pi n}\left[1-2\cos(n\pi/2)+\cos n\pi\right].
\end{align*}
括号在 $n\equiv2\pmod4$ 时为4，其他整数 $n$ 时为0。因此
\[
 \operatorname{sq}(t)\sim\sum_{j=0}^\infty\frac8{\pi(4j+2)}\sin\frac{(4j+2)t}{2}
 =\frac4\pi\sum_{j=0}^\infty\frac{\sin((2j+1)t)}{2j+1}.
\]
改用较大的周期只使频率编号加倍，实际级数不变。
(c) 平移 $t$ 得
\[
 \operatorname{sq}(t-\pi/4)\sim\frac4\pi\sum_{\substack{n\geq1\\n\text{为奇数}}}
 \frac{\cos(n\pi/4)\sin nt-\sin(n\pi/4)\cos nt}{n}.
\]
最小周期仍为 $2\pi$，跳点为 $\pi/4+k\pi$，其间高低值交替。
(d) 伸缩及竖直平移给
\[
 1+2\operatorname{sq}(2\pi t)\sim1+\frac8\pi\sum_{j=0}^\infty
 \frac{\sin(2\pi(2j+1)t)}{2j+1}.
\]
最小周期1，每周期前半值3、后半值 $-1$，跳点值1。
(e) 由 $f(t)=\frac\pi4\operatorname{sq}(t+\pi/2)$，得
\[
 f(t)\sim\sum_{j=0}^\infty\frac{(-1)^j}{2j+1}\cos((2j+1)t),
\]
最小周期 $2\pi$，与 II.20 由积分所得结果一致。
(f) $g$ 是连续奇三角波，除角点外 $g'=(4/\pi)f$，故积分后候选为
\[
 g(t)=\frac4\pi\sum_{j=0}^\infty\frac{(-1)^j}{(2j+1)^2}\sin((2j+1)t).
\]
这里可严谨地积分级数：$f$ 的部分和在 $L^2[-\pi,\pi]$ 中收敛到 $f$；对任意 $t$，积分误差的绝对值由柯西--施瓦茨不等式控制为 $\sqrt{2\pi}\,\|f-S_N\|_2$，故积分一致收敛。每个积分后的系数为 $O(n^{-2})$，级数也绝对一致收敛；常数由 $g(0)=0$ 确定。所得积分在每个开区间的导数为 $\pm1$，且连续接合，正是题给 $g$，最小周期 $2\pi$。不能在角点声称 $g'$ 存在。

(c)、(d)、(e)、(f)的函数图如下，(e)也可参照 II.20 的较大图。竖线仅为跳跃位置，(c)跳点取0、(d)跳点取1、(e)跳点取0；(f)是连续折线。
\begin{center}\begin{tikzpicture}
\begin{axis}[name=cplot,width=.43\linewidth,height=3.7cm,axis lines=middle,title={(c)},xlabel={$t$},xmin=-3.1416,xmax=3.1416,ymin=-1.3,ymax=1.3]
\addplot[blue,thick]coordinates{(-3.14159,1)(-2.35619,1)(-2.35619,-1)(.785398,-1)(.785398,1)(3.14159,1)};
\addplot[blue,only marks]coordinates{(-2.35619,0)(.785398,0)};
\end{axis}
\begin{axis}[name=dplot,at={(cplot.east)},anchor=west,xshift=1.2cm,width=.43\linewidth,height=3.7cm,axis lines=middle,title={(d)},xlabel={$t$},xmin=0,xmax=2,ymin=-1.5,ymax=3.5]
\addplot[blue,thick]coordinates{(0,3)(.5,3)(.5,-1)(1,-1)(1,3)(1.5,3)(1.5,-1)(2,-1)};
\addplot[blue,only marks]coordinates{(0,1)(.5,1)(1,1)(1.5,1)(2,1)};
\end{axis}
\begin{axis}[name=eplot,at={(cplot.south)},anchor=north,yshift=-1.2cm,width=.43\linewidth,height=3.7cm,axis lines=middle,title={(e)},xlabel={$t$},xmin=-3.1416,xmax=3.1416,ymin=-1,ymax=1]
\addplot[blue,thick]coordinates{(-3.14159,-.785398)(-1.570796,-.785398)(-1.570796,.785398)(1.570796,.785398)(1.570796,-.785398)(3.14159,-.785398)};
\addplot[blue,only marks]coordinates{(-1.570796,0)(1.570796,0)};
\end{axis}
\begin{axis}[at={(dplot.south)},anchor=north,yshift=-1.2cm,width=.43\linewidth,height=3.7cm,axis lines=middle,title={(f)},xlabel={$t$},xmin=-3.1416,xmax=3.1416,ymin=-1.9,ymax=1.9]
\addplot[blue,thick]coordinates{(-3.14159,0)(-1.570796,-1.570796)(1.570796,1.570796)(3.14159,0)};
\end{axis}\end{tikzpicture}\end{center}
\end{solution}
